\section{Baseline Evaluation}
\label{sec:baseline}

Before examining the results of the Tripartite RAFT pipeline,
it is essential to establish a baseline by evaluating the performance of LLMs in the absence of retrieval-augmented generation.
This preliminary assessment highlights the inherent limitations of LLMs when operating without external context, thereby motivating the need for the proposed approach.

\subsection{Illustrative Example}

Consider the following example. \ref{appendix:zero_shot_without_context_answer} presents the response from "Zero-shot Baseline" \ref{method:zero-shot-baseline} generated by Google Gemma 4 E4B~\cite{gemma4} to the simple question: \emph{``What is an AI Agent?''}

A critical examination of this response reveals several significant shortcomings:

\begin{itemize}
    \item \textbf{Irrelevant coverage of LLMs:} The course materials from which this question was drawn do not discuss large language models. The model's answer introduces topics that are outside the scope of the relevant curriculum.
    \item \textbf{Erroneous comparison between chatbots and AI agents:} The response conflates the concept of a chatbot with that of an AI agent, treating them as equivalent or interchangeable, which is conceptually inaccurate.
    \item \textbf{Mischaracterization of traditional AI:} The model refers to traditional AI systems as ``if-then'' rules, which is an incorrect and reductive characterization of classical artificial intelligence methodologies.
    \item \textbf{Unsubstantiated claim about memory:} The answer asserts that AI agents possess memory; however, this property is neither mentioned nor demonstrated in the course materials.
    \item \textbf{Hallucinations:} The response contains factual inaccuracies and fabricated information that cannot be verified against the source material.
\end{itemize}

These issues collectively demonstrate that LLMs, when operating without domain-specific contextual grounding, are prone to generating responses that are inaccurate, irrelevant, or misleading. This underscores the necessity of a retrieval-augmented framework such as the one proposed in this thesis.

Another example is the answer(\ref{appendix:context-augmented-baseline-answer}) of the same question from "Context-augmented Baseline" \ref{method:context-augmented-baseline} by the same model.